\section*{Resumen ejecutivo}
\addcontentsline{toc}{section}{Resumen ejecutivo}
%\begin{orientacion}
%Redacte entre 150 y 250 palabras. Presente la necesidad, el objetivo, la solución conceptual, las decisiones principales de control y comunicaciones, y el resultado alcanzado durante el segundo corte. No anuncie resultados experimentales que todavía no se hayan obtenido.
%\end{orientacion}
Los procesos de limpieza en la industria de alimentos y bebidas requieren ciclos repetibles, controlados y trazables que garanticen la inocuidad del producto. Sin embargo, muchos laboratorios y pequeñas plantas piloto no cuentan con una unidad de limpieza en sitio (Clean-In-Place, CIP) y dependen de procedimientos manuales difíciles de estandarizar \cite{bowser2023cip}. Este proyecto tiene como objetivo diseñar e implementar un prototipo de sistema portátil de limpieza CIP que prepare, caliente, dosifique e impulse la solución de limpieza a través de una tubería de prueba, en un esquema de paso único hacia un tanque de efluentes, con las fases de preenjuague, lavado con detergente y enjuague final.
En el segundo corte se construyó el prototipo sobre dos carros. Cuenta con un tanque de mezcla de 50\,L (20\,L por fase), una resistencia de inmersión de 2000\,W, un sensor PT100 con transmisor de 4--20\,mA, un flotador doble, bombas de impulsión y de agitación y una electroválvula de llenado, gobernados por un PLC S7-1200 mediante SSR y relés. Se programaron una versión preliminar de la secuencia con las tres fases y un controlador PI de temperatura. Un ensayo en lazo abierto permitió identificar el modelo térmico del tanque, con una constante de tiempo de 88,4\,min. La supervisión remota se implementó con una Raspberry Pi con Node-RED, una API en AWS y una interfaz en Streamlit, sin intervenir en el control local. Quedan para el tercer corte la instalación de la bomba dosificadora y de la electroválvula de descarga, el ajuste del controlador en lazo cerrado y la validación del prototipo frente a los requisitos.

\textbf{Palabras clave:} limpieza en sitio (CIP); sistema portátil; PLC; control de temperatura; supervisión IoT.
\section*{Control de versión del documento}
\addcontentsline{toc}{section}{Control de versión}
\begin{longtable}{|c|c|p{7cm}|c|}
\caption{Historial de versiones y control de cambios}\label{tab:versiones}\\
\hline
\textbf{Hash} & \textbf{Fecha} & \textbf{Descripción / Mensaje} & \textbf{Autor} \\ \hline
\endfirsthead
\hline
\textbf{Hash} & \textbf{Fecha} & \textbf{Descripción / Mensaje} & \textbf{Autor} \\ \hline
\endhead
\multicolumn{4}{|l|}{\textit{Segundo corte}} \\ \hline
3981d06 & 2026-10-16 & Presupuestos, diagrama unifilar y evidencias fotográficas & kvillaronga \\ \hline
fa80196 & 2026-10-16 & Unifilar en 8.1; XV-103 pendiente de instalación; TK-101 de 1,8\,L & Kare\_g \\ \hline
7310215 & 2026-10-16 & Fotos del montaje; conclusiones y trabajo previsto del segundo corte & Kare\_g \\ \hline
2399a90 & 2026-10-16 & Programa de control de temperatura y Anexo J con los programas del PLC & Kare\_g \\ \hline
f125cb9 & 2026-10-16 & Ladder v1.1 con las tres fases y GRAFCET preliminar & Kare\_g \\ \hline
a5dfbd7 & 2026-10-16 & P\&ID según ISA-5.1 e ISO 10628-2 & Kare\_g \\ \hline
15074e8 & 2026-10-16 & Secciones de montaje físico y pruebas & kvillaronga \\ \hline
17cd73c & 2026-10-16 & Validación experimental del modelo térmico (8.5) & Kare\_g \\ \hline
0022ff0 & 2026-10-16 & Sección 4: diseño conceptual ajustado al prototipo & Kare\_g \\ \hline
211890b & 2026-10-16 & Sección 6: comunicaciones de la arquitectura implementada & Kare\_g \\ \hline
54d8241 & 2026-10-16 & Sección 7: plan de validación del segundo corte & Kare\_g \\ \hline
5dc4063 & 2026-10-15 & Montaje físico: instrumentación electrónica (8.1) & Kare\_g \\ \hline
63211f1 & 2026-10-15 & Montaje físico: comunicación (8.3) y ajustes en control (8.5) & Kare\_g \\ \hline
2512a07 & 2026-10-15 & Cambios generales y numeración de tablas & kvillaronga \\ \hline
98f12aa & 2026-10-15 & Sección 2: caracterización del proceso CIP & Kare\_g \\ \hline
de3981b & 2026-10-15 & Sección 1, portada y ajustes de control & Kare\_g \\ \hline
5d17344 & 2026-10-15 & Estructura del segundo corte y registro de decisiones & Kare\_g \\ \hline
\multicolumn{4}{|l|}{\textit{Primer corte}} \\ \hline
5956a5e & 2026-09-18 & Merge branch 'main' of https://github.com/nicami13/Proyecto\_CIP\_2026-2 & fabianruizma2404-cmyk \\ \hline
618f4fa & 2026-09-18 & Revisión & fabianruizma2404-cmyk \\ \hline
bcb409b & 2026-09-18 & Cambios sobre el merge anterior & kvillaronga \\ \hline
423c8c2 & 2026-09-18 & Merge branch 'main' of https://github.com/nicami13/Proyecto\_CIP\_2026-2 & Julian Jimenez \\ \hline
481dfa3 & 2026-09-18 & Ajuste en matriz Pugh: eliminacion de elementos innecesarios, referencias de tablas, arreglo de inconsistencias respecto a la recirculacion. Actualizacion de tabla de presupuesto & Julian Jimenez \\ \hline
9f753a5 & 2026-09-18 & Cambios Karen y Joseph & kvillaronga \\ \hline
cbf13c0 & 2026-09-18 & Actualización del documento y sus correcciones & YULIETH \\ \hline
a2e2905 & 2026-09-18 & Corrección sección comunicación y anexo & FerneyPatino \\ \hline
6a745df & 2026-09-18 & Actualización presupuesto y matriz QFD & Kare\_g \\ \hline
2dd74cd & 2026-09-18 & Actualizar tabla de requisitos del proyecto & Kare\_g \\ \hline
\end{longtable}


